\section{A sharpened physical consumer and its simulation budget}
\label{sec:new-consumer}
The exact two-preparation theorem has a downstream use with explicit
hypotheses, rather than a merely renamed dependency arrow. It sharpens
the mean margin of a raw-interface confidence test, and the causal
morphism theorem transfers that test, with its full memory and
calibration costs, to another executable interface.

\begin{theorem}[Confidence and transported resource certificate]
\label{thm:transported-consumer}
Let $Q$ be a physical marked target with $\TV(Q,Q_*)\le\beta$.
Use the two-preparation event selector with a dyadic coin $\tau_b$
satisfying $|\tau_b-\tau|\le2^{-b}$, and put
\[
 a_b=v_2-\beta-\frac5{16}2^{-b},\qquad
 W_b=\max\{12,6(b+1)+5\}.
\]
If $0<c<a_b$, $R$ independent four-preparation blocks, followed by the
test $\overline D_R>c$, have error bounds
\begin{align}
 \sup_{\rm private}\Pp(\overline D_R\le c)
       &\le \exp\{-R(a_b-c)^2/2\},\label{eq:consumer-private}\\
 \sup_{\rm equal}\Pp(\overline D_R>c)
       &\le \exp\{-Rc^2/2\}.                         \label{eq:consumer-null}
\end{align}
They use $4R$ preparations, at most $bR$ fresh random bits, and at most
$W_b(2R+1)$ persistent states under a clocked schedule. With an exact
atomic coin, replace $a_b,W_b$ by $v_2-\beta,12$.

Suppose a source interface simulates this \emph{entire} repeated audit
under Theorem~\ref{thm:joint-transport}, uniformly over the private and
equal-law families, with path defect $\Delta_R$. Then the corresponding
source error bounds are the right sides of
\eqref{eq:consumer-private}--\eqref{eq:consumer-null} plus $\Delta_R$,
capped at one. At aligned cuts its state bound is
$C_{\rm cal}R_tM_tW_b(2R+1)$, and its preparation count is
$N_{\rm cal}$ plus all source calls required to supply the $4R$ audit
preparations. Calibration is acquired and retained, not supplied as a
free real-valued table.
\end{theorem}
\begin{proof}
Corollary~\ref{cor:two-physical-exact} and
Proposition~\ref{prop:tau-precision} give mean block score at least $a_b$
under every fixed private parameter. Independent fresh preparations and
independent coins give independent block scores in $[-1,1]$.
Hoeffding's inequality proves \eqref{eq:consumer-private}. Under the
equal-law null, conditional on training and the frozen event, the two
validation indicators have the same mean. Each block therefore has
mean zero, proving \eqref{eq:consumer-null} by the same inequality.
Retain the integer sum in $\{-R,\ldots,R\}$ alongside the current block
register. Updating the sum and resetting the block is a deterministic
boundary operation. This proves the raw resource bounds.

Apply Theorem~\ref{thm:joint-transport} to the whole repeated experiment,
including its stopping and validation fields. Its defect bound applies
to each of the two terminal error events, whose indicators have range
one; each probability therefore changes by at most $\Delta_R$, not
$2\Delta_R$. The same theorem gives the state product and preparation
count. The source blocks need not be independent after a shared
calibration stage: independence was used only for the reference audit,
and path-law comparison transfers its final event probability.
\end{proof}

The first part uses the exact saddle, the original microscopic bridge
and the compiled residual machine. The second additionally needs the
uniform acquired causal simulation theorem. A source which does not
satisfy that theorem's stability and calibration hypotheses is not
covered. Likewise this consumer supplies no reverse simulation and
therefore does not transfer an optimal-width lower bound to an arbitrary
source. It is an actual conditional theorem interface for subsequent
models, not an assertion that the independent geometric A2 chain or
the unbounded B4/C2 operator obligations have already been discharged.
